\documentclass[11pt]{article}
\usepackage[T1]{fontenc}
\usepackage{lmodern}
\usepackage{amsmath,amssymb,amsthm}
\usepackage[margin=1in]{geometry}
\usepackage{booktabs}
\usepackage[hidelinks]{hyperref}
\usepackage{url}
\setlength{\emergencystretch}{2em}
\newtheorem{theorem}{Theorem}
\newtheorem{lemma}[theorem]{Lemma}
\newtheorem{proposition}[theorem]{Proposition}
\newcommand{\N}{\mathbb N}
\newcommand{\E}{E_c}
\newcommand{\card}[1]{\lvert #1\rvert}
\newcommand{\ind}[1]{\mathbf{1}_{#1}}
\title{Eliahou nonnegativity at depth four\\with a two-element first interval}
\author{Zhengyi Liu}
\date{November 26, 2026}

\begin{document}
\maketitle
\begin{abstract}
Let $S$ be a numerical semigroup of multiplicity $m$ and conductor $c$.
We prove that its Eliahou number is nonnegative whenever
$3m<c\leq4m$ and $\card{S\cap[m,2m)}=2$.
There is no restriction on the distance between the two elements or on
the number of higher primitive elements.
The argument classifies short primitive representations and bounds
truncated translates of a finite set; it does not rely on enumeration of
semigroups. We prove the stronger statement for any cutoff $c$ in this
range with $[c,\infty)\subseteq S$, and provide a complete Lean 4 proof
from the semigroup assumptions to the inequality. The conclusion concerns
this depth-four family, not the general Wilf conjecture or the general
small-gap conjecture.
\end{abstract}

\section{Statement and relation to previous results}

Throughout, $\N=\{0,1,2,\ldots\}$, and intervals denote integer intervals
with the indicated endpoints. A numerical semigroup is an additive
submonoid of $\N$ with finite complement. Its multiplicity is its least
positive element. An element is \emph{primitive} if it is positive and
cannot be written as a sum of two positive semigroup elements. Write
$P$ for the primitives and $D=(S\setminus\{0\})+(S\setminus\{0\})$ for
the decomposable elements.

Fix a positive cutoff $c$ with $[c,\infty)\subseteq S$, and put
\[
 L=S\cap[0,c),\quad q=\lceil c/m\rceil,\quad \rho=qm-c,\quad
 D_c=D\cap[c,c+m).
\]
We use the cutoff-relative Eliahou number
\[
 E_c(S)=\card{P\cap L}\card{L}-q\card{D_c}+\rho.
\]
When $c$ is the least conductor, this is the usual Eliahou number
$E(S)$, also denoted $W_0(S)$ in the literature. Since $[c,c+m)$
has $m$ elements, the same definition is
$\card{P\cap L}\card{L}+q\card{P\cap[c,c+m)}-c$.

\begin{theorem}\label{thm:main}
Suppose $S$ has multiplicity $m$, contains $[c,\infty)$, and satisfies
\[
 3m<c\leq4m,\qquad S\cap[m,2m)=\{m,m+d\},\qquad 0<d<m.
\]
Then $E_c(S)\geq0$.
In particular, if $c$ is the conductor, then $E(S)\geq0$.
\end{theorem}

The interval $[m,2m)$ in this theorem is unshifted. It must not be
identified with the conductor-shifted first Ap\'ery layer when $\rho>0$.

Eliahou proved nonnegativity for depth at most three and under a
true-grading condition \cite{eliahou}. Eliahou and Mar\'in-Arag\'on
proved it when $\card L\leq12$ \cite{twelve}. Their treatment of first
Ap\'ery count one at depth four retains that cardinality restriction.
Eliahou's divset criterion yields nonnegativity for a tame model when
$m\mid c$ \cite{divsets}. Other sufficient conditions use concentration,
the maximum gap between consecutive positive semigroup elements
\cite{concentration}. These results cover parts of the present family.
Section~\ref{sec:scope} gives a family outside the hypotheses of the
specific criteria just listed. This comparison does not assert that those
criteria exhaust the literature or certify that no prior equivalent
statement exists.

Depth four admits negative Eliahou numbers: for example,
$\langle14,22,23\rangle\cup[56,\infty)$ has $E=-1$
\cite{nearmisses}. Its first interval has three elements.
Theorem~\ref{thm:main} addresses a restricted infinite class in the
direction of Bacher's conjecture that a second primitive at most $m+6$
implies nonnegativity \cite{bacher}. The theorem permits every gap $d<m$,
but requires depth four and exactly two first-interval elements.

\section{Representations and the counting reduction}

Write
\[
 c=3m+h,\qquad 1\leq h\leq m,\qquad \rho=m-h,\qquad a=m+d.
\]
Let $G=\langle m,a\rangle$ be the additive monoid generated by $m,a$;
it need not itself be cofinite. Partition the left primitives other than
$m,a$ into
\[
 B=P\cap[2m,2m+h)=2m+J,\qquad
 C=P\cap[2m+h,3m+h).
\]
Here $J\subseteq[0,h)$, and set $t=\card J$, $u=\card C$.
The indices $0,d$, and $2d$ when $2d<m$, are excluded from $J$:
the corresponding elements are decomposable.

Define
\[
 \begin{aligned}
 P_J&=(J+d)\cap[0,h),& U&=J\cup P_J,\\
 V&=\{x\in U:3m+x\notin G\},& v&=\card V,\\
 k_0&=\card{G\cap[0,c)},&
 s_0&=\card{G\cap[c,c+m)},\\
 H_J&=\{j\in J:j+d\geq h\},&
 Q_J&=(J+\{0,d,2d\})\cap[0,h),\\
 w&=\card{Q_J}+\card{H_J}.&&
 \end{aligned}
\]
All elements counted by $s_0$ are decomposable.

\begin{lemma}\label{lem:counts}
Writing $k=\card L$, $\ell=\card{P\cap L}$ and $s=\card{D_c}$, one has
\[
 k=k_0+t+v+u,\qquad \ell=2+t+u,\qquad
 s\leq s_0+w+\frac{t(t+1)}2+2u.
\]
Moreover $k_0\geq6$.
\end{lemma}
\begin{proof}
Every positive semigroup element is a sum of primitives, by induction on
its value. A representation below $c$ containing an element of $C$
can only consist of that single primitive, since adding $m$ reaches $c$.
Two elements of $B$ already sum to at least $4m\geq c$.
A representation with one element of $B$ therefore has at most one further
primitive, which must be $m$ or $a$. Consequently the following is a
disjoint union:
\[
 L=(G\cap[0,c))\ \sqcup\ B\ \sqcup\ (3m+V)\ \sqcup\ C.
\]
Primitivity and the definition of $V$ ensure the disjointness. This proves
the first two formulas.

A decomposable element of $[c,c+m)$ has a primitive representation using
only left primitives: any primitive at least $c$, plus a positive
summand, reaches $c+m$. A representation with one $b=2m+j\in B$ and no
$C$ can contribute only
\[
 b+a\quad(j+d\geq h),\qquad
 b+2m+i d\quad(j+i d<h,\ i=0,1,2).
\]
The value $b+m$ is below $c$; three additional first-interval primitives
would reach $5m\geq c+m$. The displayed contributions have cardinality
at most $w$. Two $B$ elements contribute at most $t(t+1)/2$ unordered
pairs, including repeated summands. A further positive primitive is too
large. Finally, a representation involving $C$ consists of one $C$
element and either $m$ or $a$: adding two further positive primitives
reaches $c+m$, and two $C$ elements exceed it because $c>3m$.
This contributes at most $2u$.

The six distinct elements $0,m,a,2m,m+a,3m$ belong to $G\cap[0,c)$,
proving the last assertion.
\end{proof}

Set
\[
 E_0=2k_0-4s_0+\rho,\qquad
 F=E_0+(k_0+2)t+(t+2)(v-t)-4w.
\]
Expanding the inequalities of Lemma~\ref{lem:counts} gives
\begin{equation}\label{eq:core}
 E_c(S)\geq F+u(k_0+2t+v-6)+u^2\geq F.
\end{equation}
It remains to prove $F\geq0$. This direct calculation avoids changing
the semigroup or its conductor.

\section{Truncated translates}

\begin{lemma}\label{lem:shifts}
Let $\beta=\card{P_J\setminus J}$. Then
\[
 0\leq\beta\leq t,\qquad \card U=t+\beta,
 \qquad w\leq3t,\qquad w\leq2t+2\beta.
\]
If $h\leq2d$, then $w\leq2t$.
If $P_J\subseteq J$, then also $w\leq2t$.
Finally, if $\card{U\setminus\{z\}}<t$, then
\[
 \card{U\setminus\{z\}}=t-1,\qquad z\in J,\qquad P_J\subseteq J.
\]
\end{lemma}
\begin{proof}
Translation is injective, so $\card{P_J}+\card{H_J}=t$ and
$\card{P_J}\leq t$. The second translate $(J+2d)\cap[0,h)$ has
cardinality at most $t$. Thus
\[
 \card{Q_J}+\card{H_J}
 \leq t+\card{P_J}+t+\card{H_J}=3t.
\]
If $h\leq2d$, that second translate is empty and the upper bound is $2t$.

Every element $x$ of the second translate outside $J\cup P_J$ maps
injectively to $x-d\in P_J\setminus J$. Therefore
$\card{Q_J}\leq t+2\beta$, yielding $w\leq2t+2\beta$.
If $P_J\subseteq J$, applying the truncated shift closure twice gives
$Q_J=J$, whence $w\leq2t$.

For the final statement, $J\subseteq U$ gives $\card U\geq t$.
Deleting one element can reduce the cardinality by at most one.
The strict inequality therefore forces $\card U=t$, $z\in U$, and
$U=J$. The conclusions follow.
\end{proof}

\section{Proof of Theorem~\ref{thm:main}}

We use three cases according to the position of $2d$ relative to $m$.
The finite lists below classify representations with at most four
primitive summands. They impose no bound on $m,d,h,t$, or $u$.

\subsection{The case \(2d<m\)}

The base elements below $3m$ are exactly
$0,m,a,2m,m+a,2a$. Those in $[3m,c)$ have offsets
$\{0,d,2d,3d\}\cap[0,h)$. Hence
\[
\begin{array}{c|c|c|c}
\text{range}&k_0&s_0\text{ at most}&E_0\text{ at least}\\ \hline
h\leq d&7&4&\rho-2\\
d<h\leq2d&8&4&\rho\\
2d<h\leq3d&9&4&\rho+2\\
3d<h&10&5&\rho
\end{array}
\]
For the bounds on $s_0$, the only possible threshold values are
$3m+i d$ for $i=1,2,3$ and $4m+i d$ for $i=0,1,2,3,4$.
For each $i=1,2,3$, at most one of $3m+i d,4m+i d$ lies in an interval
of length $m$. The $i=0$ value contributes one, while the $i=4$ value
can occur only if $h>4d$. Coincidences only improve the upper bounds.

The exclusions $0,d,2d\notin J$ imply that $U$ can meet the base high
offsets only at $3d$, coming from $J$ itself.
In the first two rows $v\geq t$ and $w\leq2t$, so
\[
 F\geq E_0+(k_0-6)t+(t+2)(v-t)\geq0.
\]
In the first row we used $\rho=m-h\geq m-d\geq2$.

In the third row there is no overlap and $v=t+\beta$.
Lemma~\ref{lem:shifts} gives
\[
 F\geq \rho+2+3t+(t-6)\beta.
\]
For $t\geq6$ this is nonnegative. For $0\leq t\leq5$, using
$\beta\leq t$ gives
\[
 F\geq \rho+(t-1)(t-2)\geq0,
\]
where the last inequality holds for integer $t$.

In the fourth row, $v\geq t$ together with $w\leq3t$ suffices.
Otherwise $V=U\setminus\{3d\}$ has size below $t$.
Lemma~\ref{lem:shifts} gives $v=t-1$, $3d\in J$, and $P_J\subseteq J$.
In particular $t\geq1$ and $w\leq2t$, whence
\[
 F\geq \rho+3t-2>0.
\]

\subsection{The case \(2d=m\)}

Here $2a=3m$. One has $k_0=6$ for $h\leq d$ and $k_0=7$ for $h>d$,
while $s_0=2$ in both cases.
The base high offsets are $\{0,d\}\cap[0,h)$, which $U$ avoids.
Thus $v\geq t$ and $w\leq2t$, and
\[
 F\geq 2k_0-8+\rho+(k_0-6)t\geq0.
\]
Separating this case accounts for the collision $2a=3m$.

\subsection{The case \(2d>m\)}

Put $z=2d-m$ and $r=3d-m$, so $0<z<d<r$. Below $3m$ the base
elements are $0,m,a,2m,m+a$, and their high offsets are
$\{0,z,d\}\cap[0,h)$. Consequently
\[
 k_0=6+\ind{h>z}+\ind{h>d},\qquad
 s_0=3+\ind{h>r}.
\]
For the second identity the threshold contains one of $2a,2a+m$,
one of $2m+a,3m+a$, and $4m$, together with $3a$ exactly when $h>r$.
These values are distinct.

Only $z$ can be a base overlap of $U$, and it comes from $J$.
With $e=\ind{z\in J}$, we have $v=t+\beta-e$.
As $2d>m\geq h$, Lemma~\ref{lem:shifts} gives $w\leq2t$, and hence
\begin{equation}\label{eq:large}
 F\geq E_0+(k_0-6)t+(t+2)(\beta-e).
\end{equation}
All the relevant values of $E_0$ are nonnegative.
If $\beta\geq e$, we are done. Otherwise $\beta=0,e=1$,
so $z\in J$, $P_J\subseteq J$, and $v=t-1$. The remaining bounds are
\[
\begin{array}{c|c|c|c}
\text{range}&k_0&E_0&F\text{ at least in the deficient case}\\ \hline
h\leq z&6&\rho&\text{impossible, since }z\notin J\\
z<h\leq d&7&\rho+2&\rho\\
d<h\leq r&8&\rho+4&\rho+t+2\\
r<h&8&\rho&\rho+t-2
\end{array}
\]
In the last row, $z\in J$ and $z+d=r<h$ force $r\in J$ by truncated
shift closure. These are distinct indices, so $t\geq2$.
The last bound is therefore nonnegative as well. If $r\geq m$, that
row is empty. This proves $F\geq0$ in every case; equation~\eqref{eq:core}
completes the proof of Theorem~\ref{thm:main}. \qed

\section{An explicit scope comparison}\label{sec:scope}

For every $m\geq42$, consider
\[
 S_m=\langle m,m+1,2m+10,2m+20\rangle\cup[4m-1,\infty).
\]
Its left part is exactly
\[
 \{0\}\ \cup\ (m+\{0,1\})\
 \cup\ (2m+\{0,1,2,10,20\})\
 \cup\ (3m+\{0,1,2,3,10,11,20,21\}).
\]
Indeed these are precisely the representations below $4m-1$, classified
by the number of generators involved. In particular $4m-2$ is a gap,
so the stated cutoff is the actual conductor. Thus $\card L=16$ and
$\card{P\cap L}=4$. The threshold decomposables, expressed as offsets
from $4m$, are
\[
 \{0,1,2,3,4,10,11,12,20,21,22,30,40\}.
\]
Even at $m=42$ the last offset lies below the threshold's upper endpoint.
It follows that $s=13$, $\rho=1$, and $E(S_m)=13$.

This family exceeds the twelve-left-elements hypothesis of \cite{twelve}.
Its shifted first Ap\'ery count is one, so the vanishing-early-layer
criterion in that paper does not apply. Its second shifted slice contains
the two additional primitives and is not twice its first slice; hence it
does not satisfy the true-grading equality in \cite{eliahou}.
Also $m\nmid4m-1$, outside the special-semigroup hypothesis of the cited
divset criterion \cite{divsets}.

Its concentration is $K=m-1$: this gap occurs between $m+1$ and $2m$,
and inspection of the displayed blocks shows that no larger gap occurs.
The left primitive count four satisfies neither $4\geq2K$ nor
$4\geq K+1$, required respectively in Corollaries 5.3 and 5.4 of
\cite{concentration}; the sufficient inequality in their Theorem 4.4
also fails. These observations distinguish the scope from those particular
results. They do not amount to an exhaustive originality certificate.

\section{Formal proof and reproducibility}

The full main theorem is formalized in Lean 4.36.0, with Mathlib pinned
to commit
\begin{center}
\texttt{fc4f59e577feb1a4956671bc04ac18aaaf8cd842}.
\end{center}
The checked declaration is
\begin{verbatim}
Eliahou.DepthFour.depth_four_nonneg
  (A : Admissible S m d h) : 0 <= eliahou S m h
\end{verbatim}
Here \texttt{Admissible} contains positivity and interval conditions,
zero membership, additive closure, the multiplicity lower bound,
the exact two-element first interval, and membership of the full tail
starting at $3m+h$. It contains no counting inequality or conjectural
assumption. The invariant uses integer subtraction.

The development proves the intrinsic semigroup-to-counting bridge,
the truncated-shift bounds, the unordered-pair bound, all parameter
ranges, and the final inequality. The theorem depends only on Lean's
standard \texttt{propext}, \texttt{Classical.choice}, and
\texttt{Quot.sound}; no placeholder proof or new axiom occurs in its
dependency closure. Finite enumeration of numerical semigroups and
external numerical computation are not used by this theorem.

Source, pinned dependency bootstrap, and the verification command are at
\begin{center}
\url{https://github.com/Onetkn/eliahou-depth-four-lean}.
\end{center}
After dependency setup, running \texttt{python verify.py} in the
\texttt{lean} directory builds the development and audits its axioms.
The repository's older finite examples are exploratory supplements,
separate from the uniform proof. The exact-count illustration in
Section~\ref{sec:scope} is justified by the lists given there; the main
formal declaration concerns Theorem~\ref{thm:main}.

\paragraph{AI assistance.}
OpenAI Codex assisted with the mathematical derivation, literature search,
Lean development, and manuscript drafting. Kernel checking establishes the
stated formal implication; it does not establish originality, completeness
of the literature comparison, or journal acceptance.

\begin{thebibliography}{9}
\bibitem{eliahou}
S.~Eliahou, \emph{Wilf's conjecture and Macaulay's theorem},
\href{https://arxiv.org/abs/1703.01761}{arXiv:1703.01761}, 2017.

\bibitem{twelve}
S.~Eliahou and D.~Mar\'in-Arag\'on,
\emph{On numerical semigroups with at most 12 left elements},
\href{https://arxiv.org/abs/2006.01480}{arXiv:2006.01480}, 2020.

\bibitem{nearmisses}
S.~Eliahou and J.~Fromentin, \emph{Near-misses in Wilf's conjecture},
\href{https://arxiv.org/abs/1710.03623}{arXiv:1710.03623}, 2017.

\bibitem{divsets}
S.~Eliahou, \emph{Divsets, numerical semigroups and Wilf's conjecture},
\href{https://hal.science/hal-04234167v4}{HAL-04234167v4}, 2024;
published in \emph{Communications in Algebra}, 2025,
\href{https://doi.org/10.1080/00927872.2024.2428807}{doi:10.1080/00927872.2024.2428807}.

\bibitem{concentration}
P.~Almir\'on and J.~J.~Moyano-Fern\'andez,
\emph{Eliahou number, Wilf function and concentration of a numerical semigroup},
\href{https://arxiv.org/abs/2104.03793v2}{arXiv:2104.03793v2}, 2021.

\bibitem{bacher}
A.~Bacher, \emph{Families of Eliahou semigroups linked to Farey intervals},
\href{https://arxiv.org/abs/2604.25051v9}{arXiv:2604.25051v9}, 2026.
\end{thebibliography}
\end{document}
